% ======================================================================
% SECTION 5: LIMITATIONS AND FUTURE WORK
% Two thematic blocks: per-simulation limits, and future-work direction.
% ======================================================================

[PROCESSING INSTRUCTION (Limitations block - Per-simulation limits, 4 to
6 paragraphs): Read the following files in full and synthesize the
limitations across all four simulations:
\begin{itemize}[leftmargin=*]
\item new-trial/national-24-7-trial/README.md (note the Continuity-PSL
metric is approximate, hours 56 to 83 contain approximated diagrams under
extra-hours/, and no separate code agents for local compute were
provided).
\item new-trial/final-commit/final\_error\_review.md (read this file in
full and quote at least one error category by name).
\item patient-journey/README.md (note the single-patient orientation and
the lack of a national network).
\item sponsor/final\_paper/README.md and
sponsor/final\_paper/scripts/README.md (note the 24-hour scope, the
synthetic agent decision logic, and that sponsor decisions were not
reviewed by a real pharmaceutical sponsor).
\item sponsor/final\_paper/168\_hours/README.md (note that the 7-day run
was simulated at high speed, not at wall-clock 7 days, and that
escalations and patient outcomes were generated by deterministic agent
rules rather than by clinical observation).
\item sponsor/final\_paper/168\_hours/instructions/core\_i5\_6200u\_4gb/README.md
(note that the local verification was partial - not every script
executed end-to-end on the 4 GB hardware).
\end{itemize}

For each simulation, write one short paragraph that lists the most
important limit by name. Use measured language; do not overclaim. State
explicitly that:
\begin{itemize}[leftmargin=*]
\item None of the four simulations was tested on real patients. All
patient identifiers, vital signs, and outcomes are synthetic or
illustrative. The work demonstrates a computational and software
capability, not a clinical capability.
\item Additional generations may vary in output type. Claude Code Opus
4.7 Max output is non-deterministic; a re-run of any simulation could
produce different ASCII diagrams, different agent decisions, or different
PSL trajectories. The synthesis presented in Results reflects one
realization, not a population of realizations.
\item The four simulations primarily show the computational and software
benefits that LLM-scale repository context offers over current oncology-
trial AI. They do not replace TRIPOD+AI or CREMLS reporting requirements
for any deployment-grade extension.
\end{itemize}
Cite \cite{Collins2024TRIPODAI} and \cite{ElEmam2024CREMLS}.]

[PROCESSING INSTRUCTION (Future Work block - 3 to 4 paragraphs): Define
two distinct future-work tracks and explain how they extend the
computational advantages demonstrated here into real-world applications
with real patients:

(1) Track A - Single big model performing all tasks: a future generation
of Claude Code (or an equivalent foundation-model agent) holds the entire
trial repository in context and emits all patient predictions, sponsor
decisions, and FDA-bound signals from a single inference loop. This track
is the natural extension of Simulation 4 - take the 168-hour autonomous
sponsor and connect its output stream to a real-time clinical workflow,
under TRIPOD+AI and CREMLS reporting standards.

(2) Track B - Single big model that creates smaller agents performing
smaller tasks locally: a future generation uses the 1M-token-context
model only for orchestration and routing, while delegating per-task
prediction (mortality, toxicity, response) to smaller specialized agents
that can run on-site on hardware as constrained as the Core i5-6200U.
This track is the natural extension of Simulation 1's site-level
architecture and Simulation 4's local verification - the small agents
run on-site, the big model coordinates across sites and streams
endpoints to the FDA in real time.

Discuss the trade-offs between Track A (simpler architecture, higher
inference cost per decision, single point of regulatory accountability)
and Track B (more complex architecture, lower per-site inference cost,
distributed regulatory accountability across small agents). Explicitly
state that the choice between Track A and Track B is a research question,
not a settled matter, and that both tracks should be prototyped before a
real-patient trial is designed.

Close this block by naming at least three concrete future deliverables:
(i) a TRIPOD+AI / CREMLS-compliant retrospective validation of one of
the four simulations on a real-patient cohort, (ii) a public benchmark
comparing the LLM-scale repository context approach to the supervised
baselines named in Background-A and Background-B, and (iii) an
FDA-aligned RTCT pilot submission that uses the artifacts from
Simulation 4 as inputs. Cite \cite{fda2026realtime},
\cite{Collins2024TRIPODAI}, \cite{ElEmam2024CREMLS},
\cite{kawchak\_2026\_19119939}, and \cite{kawchak\_2026\_19396256}.]
